\section{Progenitors of Our Middle Earth}

\begin{quotex}
Since men become happy by acquiring happiness, and since happiness is divinity itself, it follows that men become happy by acquiring divinity. For as men become just by acquiring integrity, and wise by acquiring wisdom, so they must in a similar way become gods by acquiring divinity. Thus everyone who is happy is a god and, although it is true that God is one by nature, still there may be many gods by participation.
\flright{\textsc{Boethius}}
\end{quotex}


Plato, Proclus \& Plotinus conceived and generated the idea of a society held together by the invisible beauty of a higher world, not dependent (strictly) upon the gods nor the accidents of Nature nor the will of man, but upon the rational soul which was immortal. This ``soul'' was actually equivalent to ``spirit'' because it was conceived not as a product of Mother Nature which was unfree by virtue of birth and controlled through the process of \emph{mimesis}, but was rational in the sense that it could reach up and out to the Good. When the ancients said ``immortal'', they predicated divinity.

\begin{quotex}
An interesting fact is that, at the beginning of Christianity, the existence of these three entities in man -- body, soul and Spirit -- was upheld. Saint Paul, for example, accepted this, Saint Augustine as well. Later it was lost through the councils and decisions of the pope and the roman church. It remained as it is known to us now: body and soul. Now it would appear that the soul is the only divine thing within man and that there is nothing else. What happened to the Spirit? It has disappeared. It is striking that it has happened this way.\footnote{\url{http://www.theforbiddenreligion.com/}}
\end{quotex}

This was the basis of Plato's war against ``poetry'' -- it was a war against the process of naked mimesis\footnote{\url{http://www.dandelion-books.com/support-files/rediscovering_plato_excerpt.pdf}}.

Whatever the cause, man had been chained to his ``matrix'', and Plato's revolution was against this cycle. Plotinus and Proclus took up this theme, which had a culmination and reflection religiously in the \emph{Therapeutae}\footnote{\url{http://en.wikipedia.org/wiki/Therapeutae}} movement in Egypt, which gave the West monasticism, and also planted seeds of this thought in the very far West, appearing again in John Scotus Erigena\footnote{\url{http://plato.stanford.edu/entries/scottus-eriugena/}}.

Boethius picked up the Platonic theme and virtually transformed it. It was his version of it, in his ``golden volume not unworthy of Tully or Plato'', which was almost all that the Middle Ages would know of ancient philosophy and which they took as a new beginning, along with Augustine, who also transformed Platonism by situating the ``Ideas'' in the mind of God. This situating of the Ideas would also mutually transform ``God'' -- God became more than Zeus, and even more than the Logos.

Victor Watts remarks in his introduction to \emph{Consolation of Philosophy:}

\begin{quotex}
Chapter 6 marks the end of Socratic dialogue and rhetorical embellishment, the end of Boethius' dependence on Plato and the advance to a higher plane of argument, viz. the exposition of the two aspects of history as Providence -- the simple unchanging plan in the mind of God -- and Fate -- the ever changing distribution in and through time of all the events God has planned in His simplicity. Boethius appears to have combined two ideas; the idea of a mutable Fate governing and revolving all things, which he read of in the treatise \emph{On Providence and Fate} by the 5\textsuperscript{th} century Neoplatonist Proclus, and the idea, already touched on at the end of Book III chapter 12, of God as the ``still turning point of the world'', and idea he found in the philosophy of Plotinus. The more the soul frees itself from corporeal things, and thus, according to both Proclus and Plotinus, from Fate, the closer it approaches the stability and simplicity of the place of rest in the centre, which according to Plotinus is God, the hinge of things, or Providence, the source of freedom and consolation for Boethius.
\end{quotex}


This is the thought, that here, Love \& Fate are one, and neither is purely itself, since both are raised to a consummated height. We are close (here) to the idea of Dante that it is Love that moves the fateful stars. This ``syncretism'' or synthesizing represented the finest philosophical or exoteric exteriorization that was possible without loss, a transmutation of all that was noblest in the Dark Ages. It knows nothing of splits and dualities, but seeks the common essence, in order to be ``carried away''. CS Lewis' extended essay \emph{The Discarded Image}, which deals with the medieval ``worldview'' devotes 15 pages to Boethius. ``Until about 200 years ago, it would have been hard to find an educated man in Europe who did not love it,'' he says. To become familiar with him, is to become ``naturalized'' to the Middle Ages. ``In his account of creation, he is closer to \emph{Timaeus} than the Scriptures''. Boethius has an intimacy of approach, a prayerful approach, which is missing in Plato (who needs a \emph{tertium quid} perpetually), but the base ideas breath the same spirit, and are not frayed genetically, but enhanced.

JAK Thompson, professor emeritus of Classics at King's College till 1959, remarks that it wasn't so much Boethius creating the Middle Ages, as Stoics and Neo-Platonism influencing the Middle Ages through him, and not in a mechanical, but a ``very fine way''. Lewis thought that Boethius' image of eternity through moving time was actually an eloquent, sophisticated improvement on the theory of Plato. In fact, Lewis went so far as to claim that Platonism had an ``unfortunate'' degree of influence on the medieval worldview, though he equivocates elsewhere in stating that there was a subtle slight disharmony between the apparently compatible positions.

Arguably, the Middle Ages had virtually nothing of Plato, yet seemed to have imbibed the ``one thing necessary'' from his breath. When Aristotle made his appearance, the debate over universals becomes possible, and we enter an ``exteriorization'' phase.

\begin{quotex}
During the first part of the Middle Ages, Platonic and neo-Platonic influences dominated philosophical thinking. ``Plato himself does not appear at all, but Platonism is everywhere,'' as Gilson has said. (Gilson 1955, p. 144.\footnote{\url{http://plato.stanford.edu/entries/medieval-philosophy/notes.html\#11}}) This situation prevailed until the recovery of Aristotle in the twelfth and thirteenth centuries. Hence, even though it is sometimes still done, it is quite wrong to think of medieval philosophy as mainly just a matter of warmed-over commentaries on Aristotle. For most of the Middle Ages by far, Aristotle was of decidedly secondary importance. This of course is not to deny that when Aristotle did come to dominate, he was \emph{very} dominant indeed and his influence was immense.\footnote{\url{http://plato.stanford.edu/entries/medieval-philosophy/\#AvailabilityOfGreekTexts}}
\end{quotex}


I do not wish to disparage Aquinas, but it seems to me that there is something very ``unmedieval'' about most of his \emph{Summa}, which does not rely on the immediate perception of God in the same way that (for example) Bonaventura does in \emph{Itinerarium Mentis ad Deum}, but rather proceeds by way of a dialectic stripped of daemonic overtones, and hence, of the ``inspiration'' which had begun to guide Plato and Socrates. We begin (here) to be concerned about meta-linguistic rules and syllogisms. This may be more a matter of how he is read, but the tendency is there.

Certainly, it is easy to quarrel with Boethius' presentation of logic, just as it is easy to quarrel with Plato's presentation of Socrates. But the Middle Ages did not die because of Progress or Necessity or even scholastic nominalism. They did not die because the West lost its ``soul'' -- Victorian English had an abundance of such. Westerners had lost their ``spirit'' which Augustine, Boethius, Paul, and Plato all breathed and shared. In fact, they had lost almost any spirit at all. It will be necessary, to change this, to breath that spirit again, which recognizes that ``the longing is not less, and the Good is not gone'', that there is ``a spirit in man which answers God'', which is not subsumed or obliterated in a sea of emotion religious or otherwise, but rather is manifested or revealed. This, of course is the medieval doctrine -- ``Grace does not eradicate, but perfects Nature''.


\flrightit{Posted on 2011-05-06 by Logres}